\documentclass[11pt]{article}
\usepackage{amsmath,amssymb,amsthm}
\usepackage[margin=1.1in]{geometry}
\newtheorem{theorem}{Theorem}
\newtheorem{lemma}[theorem]{Lemma}
\newtheorem{proposition}[theorem]{Proposition}
\newtheorem{corollary}[theorem]{Corollary}
\theoremstyle{remark}
\newtheorem{remark}[theorem]{Remark}
\newcommand{\Z}{\mathbb{Z}}
\newcommand{\N}{\mathbb{N}}
\newcommand{\PP}{\mathbb{P}}
\title{Almost all consecutive integers have coalescing $3x+1$ trajectories}
\author{Arthur M. Coll\'e \and Claude\thanks{Proofs, computations and the Lean development were produced by Claude (Anthropic) working with A.~M.~Coll\'e, who directed the work. The Lean statements have not yet had an independent human review.}}
\date{2026-11-26}
\begin{document}
\maketitle

\begin{abstract}
Run the orbits of $n$ and $n+1$ under the shortcut Collatz map side by side for $k$ steps. We prove
that the proportion of $n$ whose two orbits have not yet merged is of exact order $1/\sqrt k$: it
lies between $c_1/\sqrt k$ and $c_2/\sqrt k$ for every $k\ge1$, and the same holds for $n$ and
$n+m$. In particular almost all consecutive integers have coalescing trajectories; this is Gao's
1993 conjecture, first proved by O.~J.~Said Duran in 2026, and the rate answers a question left
open there. Every statement is checked in Lean~4. Nothing here bears on the Collatz conjecture
itself.
\end{abstract}

\section{Statement}
Let $T(x)=x/2$ for even $x$ and $T(x)=(3x+1)/2$ for odd $x$, and let $C(x)=x/2$, resp.\ $3x+1$.
For $k\ge0$ let
\[
C_k=\#\{\,0\le r<2^k:\ T^k(r+2^kq)=T^k(r+1+2^kq)\ \text{for all } q\ge0\,\}.
\]
The ratio $C_k/2^k$ is nondecreasing, and it is the natural density of
$\{n: T^k(n)=T^k(n+1)\}$.

\begin{theorem}\label{main}
$C_k/2^k\to1$ as $k\to\infty$. More generally, for every $m\ge1$ the proportion of residues $r$
modulo $2^k$ with $T^k(r+2^kq)=T^k(r+m+2^kq)$ for all $q\ge 0$ tends to $1$.
\end{theorem}

\begin{corollary}[Gao's conjecture]\label{gao}
The set $\{n:\ C^{j}(n)=C^{j}(n+1)\text{ for some } j\le\log_2 n\}$ has natural density one.
For every $L$, the set of $n$ for which $n,n+1,\dots,n+L$ have a common iterate at a common time
has natural density one; for such $n$ these integers have equal total stopping times and equal
heights whenever one of them reaches $1$.
\end{corollary}

G.-G.~Gao (\emph{Discrete Math.}\ 112 (1993), 261--267) conjectured the first statement and proved
that the density is at least $0.389$, as recorded in Lagarias's annotated bibliography
(arXiv:math/0309224, entry~68); OEIS A398792 records $C_{32}/2^{32}\approx0.4357$. The conjecture
was first settled in the preprint of Said Duran cited above. We did not read Gao's paper. As far as
that preprint's abstract shows, it treats the offset $m=1$; its Remark~3.7 and Section~7 say that its proof
yields no useful rate. Whether general offsets or the values in Section~\ref{sec:checks} appear
in its body we have not checked.

\begin{theorem}[Rate]\label{rate}
There are constants $0<c_1\le c_2$ such that for all $k\ge1$
\[
\frac{c_1}{\sqrt k}\ \le\ 1-\frac{C_k}{2^k}\ \le\ \frac{c_2}{\sqrt k},
\]
and one may take $c_1=\tfrac18$ (more precisely $1-C_k/2^k\ge 1/(4\sqrt{2k+2})$). The same two-sided
bound holds, with constants depending on $m$, for the pair $n,n+m$ with any $m\ne0$.
\end{theorem}
The preprint of Said Duran cited above states that the true order of this quantity is not known
and lists a rate as its first open question; its author's public research notes
(\texttt{gaopendientes.md} in \texttt{github.com/theswordofthemorning/conjeturadecollazpython})
record the order $1/\sqrt t$ as a heuristic from random-walk return times, ``not measured''.
Theorem~\ref{rate} proves that order. We found no earlier proof (checked: that preprint, version~1;
its two companion preprints; the author's two public repositories as of 23 November 2026; web
searches), but have not seen the version under review at Qeios. The constant
$c_2$ produced by the proof is not explicit in any useful sense; numerically
$\sqrt k\,(1-C_k/2^k)\approx 11$ for $10^4\le k\le10^8$.

\paragraph{Prior work and status.} Density one is not new. O.~J.~Said Duran, \emph{Consecutive integers
coalesce with density one under the Collatz map: a formally verified proof of a conjecture of Gao}
(Zenodo, 27 September 2026, doi:10.5281/zenodo.23003524; Lean sources at
\texttt{github.com/theswordofthemorning/gao-palomar}) proves Corollary~\ref{gao} for $n,n+1$, by
what its abstract describes as a Markov chain on pairs $\{v,3^kv+c\}$ driven by fair coins with a
contraction of $\sqrt{|c|/3^k}$. The present note was written without knowledge of that preprint
and found it only afterwards; it is an independent proof by the same basic device, and we have
read only that preprint's abstract and repository page, not its proof. Theorem~\ref{main},
Corollary~\ref{gao} and Theorem~\ref{rate} below are checked by the Lean kernel
(Section~\ref{sec:checks}). Theorem~\ref{rate} appears to be new: it answers a question that
preprint leaves open (we have seen its version~1 only). Nothing here bears on the Collatz conjecture itself.

\section{The pair chain}
Extend $T$ to the $2$-adic integers $\Z_2$ with Haar measure. For $x\in\Z_2$ put
$y_k=T^k(x)$, $z_k=T^k(x+m)$, and let $a_k,a'_k$ count the odd terms among $y_0,\dots,y_{k-1}$ and
$z_0,\dots,z_{k-1}$. On the cylinder $x\equiv r \pmod{2^k}$ both are affine in $x$ with slopes
$3^{a_k}/2^k$ and $3^{a'_k}/2^k$.

Order the pair as $(p_k,q_k)$ so that $q_k$ has at least as many odd steps, let $d_k=|a'_k-a_k|$
and $\mu_k=q_k-3^{d_k}p_k$; when $d_k=0$ choose the order with $\mu_k\ge0$. Then $\mu_k\in\Z$
depends only on $r$, and the \emph{state} $X_k=(d_k,\mu_k)$ starts at $X_0=(0,m)$.

\begin{lemma}[One step]\label{step}
Let $c=p_k \bmod 2$. Then $X_{k+1}$ is obtained from $X_k=(d,\mu)$ by
\[
\begin{array}{lll}
\mu\ \text{even:} & c=0:\ (d,\ \mu/2), & c=1:\ (d,\ (3\mu+1-3^d)/2),\\
\mu\ \text{odd},\ d\ge1: & c=0:\ (d+1,\ (3\mu+1)/2), & c=1:\ (d-1,\ (\mu-3^{d-1})/2),\\
\mu\ \text{odd},\ d=0: & c=0:\ (1,\ (3\mu+1)/2), & c=1:\ (1,\ (1-3\mu)/2),
\end{array}
\]
followed by replacing $(0,\nu)$ with $(0,|\nu|)$. The coin $c$ is uniform and independent of
$x \bmod 2^k$. The state $(0,0)$ is absorbing. On a cylinder, $y_k=z_k$ identically if and only if
$X_k=(0,0)$; otherwise $y_k=z_k$ for at most one point of the cylinder.
\end{lemma}
\begin{proof}
If $\mu$ is even, $p$ and $q=3^dp+\mu$ have equal parity and the same branch of $T$ applies:
$q/2=3^d(p/2)+\mu/2$ and $(3q+1)/2=3^d\,(3p+1)/2+(3\mu+1-3^d)/2$. If $\mu$ is odd the parities
differ. For $p$ even, $(3q+1)/2=3^{d+1}(p/2)+(3\mu+1)/2$. For $p$ odd and $d\ge1$,
$q/2=3^{d-1}(3p+1)/2+(\mu-3^{d-1})/2$. For $p$ odd and $d=0$ the roles swap:
$(3p+1)/2=3\,(q/2)+(1-3\mu)/2$. The parity of $p_k$ is uniform given $x\bmod 2^k$ because $p_k$ is
affine in $t$, where $x=r+2^kt$, with odd slope. Finally $q_k-p_k=(3^d-1)p_k+\mu$ vanishes
identically iff $d=0$ and $\mu=0$, and otherwise for at most one value of $p_k$.
\end{proof}

Hence $C_k$ is the number of coin words of length $k$ leading from $(0,1)$ to $(0,0)$, the two
orbits have equally many odd steps on merged classes, and Theorem~\ref{main} is the statement that
the chain is absorbed with probability one from $(0,m)$.

\section{Every offset can merge}
\begin{lemma}\label{offsets}
From every state $(0,m)$, $m\ge1$, some coin word leads to $(0,0)$.
\end{lemma}
\begin{proof}
Strong induction on $m$. If $m$ is even, coin $0$ leads to $(0,m/2)$. If $m$ is odd write
$3m+1=2^js$ with $s$ odd. The word $0^j1$ gives
$(0,m)\to(1,2^{j-1}s)\to\dots\to(1,s)\to(0,(s-1)/2)$, and $(s-1)/2\le(3m-1)/4<m$. If $s=1$ this is
$(0,0)$.
\end{proof}
In terms of integers: for every $m\ge1$ there are $k,r$ with $T^k(r+2^kq)=T^k(r+m+2^kq)$ for all
$q\ge0$.

\section{A potential}
Write $w(d,\mu)=|\mu|/3^d$, $a_0=\tfrac12$, $a_1=\tfrac32$, and $X\cdot c$ for the state after coin
$c$.
\begin{lemma}[Domination]\label{dom}
$w(X\cdot c)\le a_c\,w(X)+\tfrac12$. If $X=(0,m)$ with $m$ odd then $w(X\cdot c)\le m/2+\tfrac16$
for both coins.
\end{lemma}
\begin{proof} Divide the formulas of Lemma~\ref{step} by the new power of $3$. \end{proof}

For $d\ge1$ let $e_d=v_2(3^d-1)$, so $1\le e_d<2d$, and
\[
\ell(d,\mu)=\#\{1\le j\le e_d:\ 2^j\mid\mu\}+[\,2^{e_d}\mid\mu\,],\qquad \ell(0,\mu)=0.
\]
\begin{lemma}\label{ell}
For $d\ge1$: $0\le\ell(d,\mu)\le e_d+1\le2d$, $\ell(d,\mu)=0$ iff $\mu$ is odd, and for $\mu$ even
\[
\ell(d,\mu/2)+\ell\bigl(d,(3\mu+1-3^d)/2\bigr)\le 2\ell(d,\mu)-2 .
\]
\end{lemma}
\begin{proof}
Put $\nu=(3\mu+1-3^d)/2$. For $j<e_d$: $2^j\mid\mu/2$ and $2^j\mid\nu$ are both equivalent to
$2^{j+1}\mid\mu$, since $2^{j+1}\mid 3^d-1$. The claim reduces to
$[2^{e_d+1}\mid\mu]+[2^{e_d}\mid\nu]\le[2^{e_d}\mid\mu]$. If $2^{e_d}\nmid\mu$ then
$2^{e_d}\nmid 3\mu-(3^d-1)$, so both terms on the left vanish. If $2^{e_d}\mid\mu$ the two terms
cannot both be $1$, as that would give $2^{e_d+1}\mid 3^d-1$.
\end{proof}

Fix $\theta=\tfrac{19}{20}$, $\lambda=\tfrac{39}{40}$, $A=250000$, $M=4\cdot10^8$ and
\[
V(d,\mu)=\sqrt d+A(1-\lambda^d)+\frac{\theta^d}{\sqrt2}\,\ell(d,\mu)+\theta^d\sqrt{w(d,\mu)},\qquad
F=\{(0,m):|m|\le M\}.
\]
\begin{proposition}\label{super}
For every state $X\notin F$: $\ \tfrac12\bigl(V(X\cdot0)+V(X\cdot1)\bigr)\le V(X)$.
\end{proposition}
\begin{proof}
Let $\rho=\tfrac12(\sqrt{1/2}+\sqrt{3/2})<1$. By Lemma~\ref{dom},
$\sqrt{w(X\cdot c)}\le\sqrt{a_c}\sqrt{w(X)}+1/\sqrt2$.

\emph{(i) $d\ge1$, $\mu$ even.} The level is unchanged. The last term of $V$ averages to at most
$\rho\,\theta^d\sqrt w+\theta^d/\sqrt2$, and by Lemma~\ref{ell} the third term drops by at least
$\theta^d/\sqrt2$.

\emph{(ii) $d\ge1$, $\mu$ odd.} The children are at levels $d\pm1$ and $\ell(X)=0$. By concavity
$\tfrac12(\sqrt{d+1}+\sqrt{d-1})\le\sqrt d$. The second term drops by exactly
$A\lambda^{d-1}(1-\lambda)^2/2=A\lambda^{d-1}/3200$. The last term averages to at most
\[
\theta^d\sqrt w\cdot\tfrac12\bigl(\theta/\sqrt2+\sqrt{3/2}/\theta\bigr)+\frac{\theta^{d+1}+\theta^{d-1}}{2\sqrt2}
\le\theta^d\sqrt w+\frac{\theta^{d-1}}{\sqrt2},
\]
because $\tfrac12(\theta/\sqrt2+\sqrt{3/2}/\theta)<0.981$. The third term averages to at most
$\frac{1}{2\sqrt2}\bigl(2(d+1)\theta^{d+1}+2(d-1)\theta^{d-1}\bigr)\le 2d\,\theta^{d-1}/\sqrt2$.
So it suffices that $(1+2d)\theta^{d-1}/\sqrt2\le A\lambda^{d-1}/3200$. By Bernoulli,
$(\lambda/\theta)^{d-1}=(39/38)^{d-1}\ge1+(d-1)/38$, whence $(1+2d)\theta^{d-1}\le76\lambda^{d-1}$,
and $76/\sqrt2<78.125=A/3200$.

\emph{(iii) $d=0$, $\mu$ even, $m=|\mu|>M$.} The children are $(0,m/2)$ and $(0,3m/2)$, and
$V$ averages to $\rho\sqrt m\le\sqrt m$.

\emph{(iv) $d=0$, $\mu$ odd, $m=|\mu|>M$.} Both children are at level $1$ with $w\le m/2+\tfrac16$ and
$\ell\le2$, so each has
$V\le1+A/40+\sqrt2\,\theta+\theta(\sqrt{m/2}+1/\sqrt6)$. This is at most $\sqrt m$ because
$\sqrt m\,(1-\theta/\sqrt2)\ge20000\cdot0.328>6253>1+A/40+\sqrt2\,\theta+\theta/\sqrt6$.
\end{proof}

\section{Proof of Theorem~\ref{main}}
Let $\PP_X$ refer to the chain started at $X$.

\begin{lemma}\label{hitF}
For every $X$: $\PP_X(\text{the chain visits }F)\ge\tfrac12$.
\end{lemma}
\begin{proof}
Let $X\notin F$, $\Lambda=4V(X)+1$ and $G=\{Y\notin F:\ V(Y)<\Lambda\}$. By
Proposition~\ref{super}, $V$ stopped on leaving $G$ is a nonnegative supermartingale, so the
probability of reaching $\{V\ge\Lambda\}$ before $F$ is at most $V(X)/\Lambda<\tfrac14$. It remains
to see that the chain cannot stay in $G$ forever. There is $L=L(\Lambda)$ such that from every
$Y\in G$ some word of length at most $L$ leaves $G$: at level $0$ use coin $0$ while $\mu$ is even
(at most $2\log_2\Lambda$ times, since $\sqrt\mu<\Lambda$) and then any coin; at level $d\ge1$
with $\mu$ even use the coin that lowers $\ell$ (Lemma~\ref{ell}; at most $2d$ times) and then,
$\mu$ being odd, coin $0$, which raises the level. Once $d\ge\Lambda^2$ we have
$V\ge\sqrt d\ge\Lambda$. Hence the probability of remaining in $G$ for $nL$ steps is at most
$(1-2^{-L})^n\to0$.
\end{proof}

\begin{proof}[Proof of Theorem~\ref{main}]
Let $u(X)$ be the probability that the chain from $X$ is never absorbed. By Lemma~\ref{offsets}
each of the finitely many states of $F$ is absorbed with positive probability; let $c_0>0$ be the
minimum. By Lemma~\ref{hitF} and the strong Markov property, $u(X)\le1-c_0/2$ for \emph{every}
state $X$. Let $U=\sup_Xu(X)\le1-c_0/2<1$. The Markov property at time $j$ gives
$u(X)=\mathbb E_X[u(X_j)]\le U\,\PP_X(X_j\ne(0,0))$, and letting $j\to\infty$,
$u(X)\le U\,u(X)$. Therefore $u\equiv0$.
\end{proof}

\begin{proof}[Proof of Corollary~\ref{gao}]
Given $\varepsilon$ choose $K$ with $C_K\ge(1-\varepsilon)2^K$. If $n$ lies in a merged class
modulo $2^K$, then $T^K(n)=T^K(n+1)$ with the same number $a\le K$ of odd steps, so
$C^{K+a}(n)=C^{K+a}(n+1)$ and $K+a\le2K\le\log_2n$ once $n\ge4^K$. The lower density is therefore
at least $1-\varepsilon$. For runs apply Theorem~\ref{main} to $m=1,\dots,L$ and intersect.
\end{proof}

\section{The rate}
\begin{proof}[Proof of Theorem~\ref{rate}, lower bound]
Absorption happens only by coin $1$ at the state $(1,1)$. The words beginning $1,0$ lead
$(0,1)\to(1,-1)\to(2,-1)$. From then on the level moves by $\pm1$ with fair independent coins at
the steps where $\mu$ is odd, and is constant otherwise. If the first $k$ of these moves keep the
level at least $2$, the chain is not absorbed by time $k$. A simple random walk stays nonnegative
for $k$ steps with probability $\binom{k}{\lfloor k/2\rfloor}2^{-k}\ge1/\sqrt{2k+2}$, and the first two coins cost a factor $\tfrac14$.
\end{proof}

For the upper bound let $\Gamma_2=V-\sqrt d$ and $\Gamma=d+\Gamma_2$, with $V$, $F$ as in
Section~4, and let $\tau_F$ be the first time the chain is in $F$.

\begin{lemma}\label{gam}
$\Gamma_2\ge0$. Off level $0$, $\Gamma_2$ is superharmonic; at level $0$ outside $F$ its average
over the two children is at most $\Gamma_2-1$, and at most $\Gamma_2-24$ when $\mu$ is odd. Hence
$\Gamma$ is superharmonic outside $F$.
\end{lemma}
\begin{proof}
These are the computations of Proposition~\ref{super} without the concavity step; the slack in
cases (iii) and (iv) there is at least $680$ and $300$ respectively.
\end{proof}

\begin{lemma}[Time below level $L$]\label{ql}
Put $\hat e_0=0$, $\hat e_j=e_j$ for $j\ge1$, $E_2(n)=\sum_{j<n}(n-j)\hat e_j$, and for $L\ge1$
\[
q_L(d)=(8L+14)d+6-2d^2-E_2(d+1)-E_2(d-1),\qquad Q_L=q_L(d)+\ell+L\,\Gamma_2 .
\]
Then $Q_L\ge0$ on levels $d\le L$, $Q_L\le(8L+16)d+6+L\Gamma_2$, and for every state outside $F$
of level $d<L$ the average of $Q_L$ over the two children is at most $Q_L-1$. Consequently the
expected time before the chain either visits $F$ or reaches level $L$ is at most $Q_L$.
\end{lemma}
\begin{proof}
$E_2$ has second difference $\hat e$, so $q_L(d+1)+q_L(d-1)-2q_L(d)=-4-\hat e_{d+1}-\hat e_{d-1}$.
Since $e_j=1$ for odd $j$ and $e_j=2+v_2(j)$ for even $j$, $\sum_{j\le d}\hat e_j\le3d$, whence
$E_2(d+1)\le3d(d+1)$ and $q_L(d)\ge6$ for $d\le L$. At a step with $\mu$ even above level $0$,
$\ell$ drops by one on average (Lemma~\ref{ell}) and nothing else rises. At a step with $\mu$ odd
above level $0$, $\ell$ rises from $0$ to at most $\tfrac12(\hat e_{d+1}+\hat e_{d-1})+1$ on
average, which the second difference of $q_L$ pays for with $1$ to spare. At level $0$ the term
$L\Gamma_2$ pays, by Lemma~\ref{gam}: $L\ge1$ when $\mu$ is even, and $24L\ge q_L(1)-q_L(0)+3$
when $\mu$ is odd.
\end{proof}

\begin{lemma}[Hitting time of $F$]\label{tail}
For every state $X$ and $t\ge1$:
$\ \PP_X(\tau_F>t)\le\bigl(33d+6+3\Gamma_2(X)\bigr)/\sqrt t$.
\end{lemma}
\begin{proof}
Let $L=\lceil\sqrt t\rceil$. By the maximal inequality for $\Gamma\ge d$ (Lemma~\ref{gam}) the
chain reaches level $L$ before $F$ with probability at most $\Gamma(X)/L$. Otherwise it stays below
level $L$ and outside $F$ for $t$ steps, which by Lemma~\ref{ql} and Markov's inequality has
probability at most $Q_L(X)/t$. Use $\sqrt t\le L\le2\sqrt t$.
\end{proof}

\begin{proof}[Proof of Theorem~\ref{rate}, upper bound]
There are $K_0$ and $c_0>0$ such that from every state of $F$ the chain is absorbed within $K_0$
steps with probability at least $c_0$ (Lemma~\ref{offsets}; $F$ is finite). Let $R$ be the finite
set of states reachable from $F$ in exactly $K_0$ steps, $G_0$ the maximum over $R$ of the
numerator in Lemma~\ref{tail}, and $H(t)$ the maximum over $R$ of the probability of not being
absorbed by time $t$. Splitting at $\tau_F$ and using the Markov property $K_0$ steps later,
\[
H(s+K_0+r)\ \le\ \frac{G_0}{\sqrt s}+(1-c_0)\,H(r).
\]
With $s=\lceil c_0t/2\rceil$ and $r=t-s-K_0\ge(1-c_0)t$ for large $t$, induction on $t$ gives
$H(t)\le B/\sqrt t$ as soon as $B(1-\sqrt{1-c_0})\ge G_0\sqrt{2/c_0}$. The start $(0,1)$ lies in
$F$, so one more application of the Markov property finishes.
\end{proof}

\begin{remark}
The argument uses the contraction $\tfrac12(\log\tfrac12+\log\tfrac32)<0$ through $\rho<1$ and the
existence of $\theta<1$, and the arithmetic of the $+1$ only through Lemmas~\ref{offsets}
and~\ref{ell}. It gives the same conclusion for $(3x-1)/2$, which is conjugate by $x\mapsto-x$. For
$(5x+1)/2$ the corresponding recursion expands on average and the method does not apply.
\end{remark}

\section{Exact values and checks}\label{sec:checks}
The chain turns $C_k$ into a transfer computation with about $1.59^k$ states instead of $2^k$
residues. A forward pass to depth $42$ joined with a backward pass of depth $24$ gives every $C_k$
for $k\le66$ (\texttt{compute/pairchain.cpp}, about one minute and $20$\,GB).
\[
\begin{array}{rrl}
k & C_k & C_k/2^k\\\hline
25 & 14030369 & 0.418138\\
32 & 1871137623 & 0.435658\\
40 & 496587379326 & 0.451644\\
50 & 526795653639168 & 0.467889\\
60 & 555030220058374659 & 0.481412\\
66 & 36051385938701759737 & 0.488587
\end{array}
\]
The values agree with the project's earlier exhaustive ledger for $k\le25$, with the $2^{32}$
figure on OEIS, with a Python implementation of the backward pass for $k\le30$, and across three
machines and two architectures using different split depths ($42$ and $39$) for $k\le61$. Monte
Carlo runs of the chain with exact arithmetic (\texttt{mc\_pairchain.py}) left $1.107\%\pm0.014\%$
of $6\cdot10^5$ paths unabsorbed at $k=10^6$ and $0.13\%\pm0.02\%$ of $4\cdot10^4$ paths at $k=10^8$;
the unabsorbed fraction is close to $11/\sqrt k$ from $k=10^4$ on, the order of
Theorem~\ref{rate}.

\medskip\noindent\textbf{Lean.} \texttt{CollatzCoalescence/} (Lean 4.33.0-rc2, Mathlib
\texttt{51e6992}; about 4700 lines; no \texttt{sorry}, no added axioms, no \texttt{native\_decide}).
\texttt{lake build} succeeds, and the concatenated sources were replayed on a second machine
against a separately built Mathlib. \texttt{\#print axioms} reports only \texttt{propext},
\texttt{Classical.choice}, \texttt{Quot.sound} for:
\begin{itemize}\itemsep0pt
\item \texttt{gao\_conjecture}: for every $\varepsilon>0$ and all large $N$, at least
$(1-\varepsilon)N$ of the $n<N$ have $C^{j}(n)=C^{j}(n+1)$ for some $j$ with $2^j\le n$
(\texttt{Gao.lean}; \texttt{coalescence\_density\_one\_log} is the same for $n,n+m$);
\item \texttt{coalescence\_density\_one}, \texttt{offset\_classes},
\texttt{merged\_proportion\_tendsto}: Theorem~\ref{main} (\texttt{Main.lean});
\item \texttt{every\_offset\_merges}, \texttt{offset\_absorbable}: Lemma~\ref{offsets};
\texttt{step\_rel}, \texttt{card\_merged\_succ}: Lemma~\ref{step} and the count;
\texttt{V\_superharmonic}: Proposition~\ref{super}; \texttt{escape} and
\texttt{G2\_lyapunov\_hitting}: Lemma~\ref{hitF}.
\end{itemize}
\texttt{neighbour\_merge\_rate} (\texttt{RateMain.lean}) is Theorem~\ref{rate}: constants
$0<c_1,c_2$ with $c_1 2^k/\sqrt k\le 2^k-C_k\le c_2 2^k/\sqrt k$ for all $k\ge1$; its inputs are
\texttt{survival\_lower}, \texttt{Gam\_super}, \texttt{QL\_step}, \texttt{one\_sub\_hit\_le} and
\texttt{survival\_rate\_of\_renewal}. \texttt{unmerged\_density\_rate} (\texttt{RateDensity.lean})
restates it for integers: for each $k\ge1$ the set of $n$ with $T^k(n)\ne T^k(n+1)$ has lower and
upper natural density between $c_1/\sqrt k$ and $c_2/\sqrt k$. \texttt{offset\_merge\_rate} (\texttt{RateOffsets.lean}) is the
same two-sided bound for $n,n+m$, $m\ne0$; its upper bound holds from every state of the chain.
The statements were read by their author only; no independent statement review and no second
proof checker has been run. Before formalization,
Proposition~\ref{super} was also checked numerically on $2.4\cdot10^5$ states.
\end{document}
